\chapter{Literature Review}
\label{sec:litreview}
\lhead{\emph{Chapter 2. Literature Review}}  % page header

This chapter reviews research about research that bears on the three research questions. Section~\ref{sec:lit-fashion}
covers fashions, herding and bandwagons in science and organisations. Section~\ref{sec:lit-novelty} covers how science
balances tradition and novelty, and how novel work is rewarded. Section~\ref{sec:lit-entry} covers newcomers and
researchers who move between topics. Section~\ref{sec:lit-review} covers peer review at ML venues,
Section~\ref{sec:lit-ml} covers meta-research on ML and natural language processing (NLP), and
Section~\ref{sec:lit-topics} covers how research topics are measured over time. Section~\ref{sec:lit-gap} summarises
what is missing.

\section{Fashions, herding and bandwagons}
\label{sec:lit-fashion}

The idea that science has fashions goes back at least to \citet{crane1969fashion}. In organisation studies,
\citet{abrahamson1991managerial} described management fashions, in which ideas spread because others adopt them rather
than because they work, and \citet{abrahamson1993institutional} showed that bandwagons can spread even when adopting
brings no gain. \citet{abrahamson1999management} followed one such fashion over its whole life and found that the way
people wrote about it differed between its rise and its fall. Institutional theory gives a related reason for
convergence: under uncertainty, organisations copy one another \citep{dimaggio1983iron}, especially organisations that
look successful \citep{haunschild1997modes}. \citet{strang2001search} modelled firms that imitate their most
successful peers and showed that this can produce cycles of adoption and abandonment.

Economic models explain why imitation can be rational for each person and still harmful overall. In an information
cascade, people follow those who chose before them and ignore their own information
\citep{bikhchandani1992theory,banerjee1992simple}. The resulting consensus can rest on very little evidence and can
break easily. Herding also protects a reputation when people are judged against their peers
\citep{scharfstein1990herd}. Diffusion research describes how a new practice spreads through a population over time,
from early adopters to late ones \citep{rogers2003diffusion}, and shows that less connected people tend to adopt later
\citep{coleman1957diffusion}.

Several authors apply these ideas to science. \citet{fujimura1988molecular} described a bandwagon in cancer research
that formed around a package of theory and methods that many laboratories could adopt. \citet{baddeley2015herding}
reviewed herding and social influence in scientific research. In formal models, \citet{weisberg2009epistemic}
separate scientists who follow others' approaches from those who ignore them, and \citet{weisberg2013modeling} model
herding as a way to avoid risk. \citet{alexander2015epistemic} dispute some of the conclusions of the first model, and
\citet{kitcher1990division} and \citet{strevens2003role} ask how credit and the priority rule shape how scientists
divide their effort. A randomised trial in ML review found no herding among reviewers in discussions
\citep{stelmakh2023large}, while a study of ICLR rebuttals found that reviewers respond to their co-reviewers' scores
\citep{kargaran2025insights}.

There is also a body of work on staying close to one's own past, which this report calls continuity. Organisations
learn from their own experience \citep{levitt1988organizational} and search near what they already know
\citep{march1991exploration,stuart1996local}. Taking up new knowledge depends on related knowledge that one already
has \citep{cohen1990absorptive}. Small early events can lock a population into one technology
\citep{arthur1989competing}. These accounts also predict that a field converges, but through each group's own record
rather than through copying others.

Two limits of this literature matter here. It mostly measures how much an idea is adopted or discussed, and rarely
observes how each adopter is judged. It also rarely separates copying others from staying close to one's own record,
although both lead to convergence.

\section{Tradition, novelty and their reward}
\label{sec:lit-novelty}

The science of science studies research itself with large data \citep{fortunato2018science}. A central theme is the
tension between building on a tradition and breaking from it \citep{kuhn1977essential}. \citet{foster2015tradition}
found that most biomedical researchers choose strategies that extend existing knowledge, and that more innovative
strategies are rarer, riskier and more variable in their payoff. \citet{uzzi2013atypical} found that the most cited
papers combine mostly conventional prior work with a small amount of unusual combinations. \citet{wang2017bias} found
that highly novel papers are less cited in the short run and more cited in the long run. At the scale of whole fields,
\citet{park2023papers} report that papers have become less disruptive over time, \citet{wu2019large} report that large
teams develop existing work while small teams disrupt it, and \citet{chu2021slowed} report that in very large fields new
papers find it harder to become widely cited. \citet{rzhetsky2015choosing} show that choices that are safe for each
scientist can make the whole field explore less. \citet{xiang2025research} review how researchers choose topics.

Whether evaluators penalise novelty is contested. In grant review, \citet{boudreau2016looking} found that evaluators
gave lower scores to highly novel proposals and to proposals close to their own expertise. In review files from 49
journals, \citet{teplitskiy2022is} found no such penalty. \citet{lane2022conservatism} showed in field experiments that
evaluators lower their scores more after seeing critical peer scores than they raise them after seeing favourable ones.
\citet{li2017expertise} found that grant reviewers favour applications in their own area, but are also better informed
about them. Cumulative advantage, in which those already recognised gain more recognition, is a further reason why a
record can matter at the gate \citep{merton1968matthew}. Priority and timing also matter to scientists
\citep{merton1957priorities}.

These studies measure reward mostly through citations, grants or published papers. They seldom see, in the same
population, both what researchers choose to write and how the gate judges every submission, including rejected ones.

\section{Newcomers and moving between topics}
\label{sec:lit-entry}

A long line of work finds that people at the margin of a field bring new ideas. On an open innovation platform,
problems were more often solved by people whose expertise lay far from the problem's field
\citep{jeppesen2010marginality}. Students from groups under-represented in their field introduce more novel ideas, but
these ideas are taken up less \citep{hofstra2020diversity}. Younger researchers try newer ideas
\citep{packalen2019age}, and fresh teams produce more novel work \citep{zeng2021fresh}. When an eminent scientist dies,
researchers from outside the field enter it and their papers are often highly cited \citep{azoulay2019does}. Inventors
who move into new fields produce more novel but less valuable patents, and the loss is smaller when they work with
experts \citep{arts2018paradise}.

Moving away from one's own past work has costs. In the pivot penalty, larger shifts from a researcher's earlier work go
with lower impact and lower publication success \citep{hill2025pivot}. Researchers who stayed close to their own record
during the deep learning shift had lower later impact \citep{bao2025persistence}. Crowding can also push researchers
away: after the Soviet collapse, American mathematicians moved out of the fields that received many Soviet
mathematicians \citep{borjas2015cognitive}. Ties to insiders help newcomers: junior scientists
who publish with established authors do better later \citep{sekara2018chaperone}, and senior expertise offsets the
disadvantage of newcomers to a new research frontier \citep{biliotti2024breaking}. In ML, half of the first
authors of LLM papers on arXiv in 2023 were new to LLM research \citep{movva2024topics}, NLP researchers increasingly publish at
ML venues \citep{jurgens2026future}, and at ICLR some papers act as catalysts that come before a topic grows
\citep{huang2026catalyst}.

This work measures how novel newcomers' work is, or how they fare later. It does not ask what type of contribution
newcomers bring, for example a new method or a benchmark of existing ones, and it rarely looks at the accept or reject
decision under double-blind review.

\section{Peer review at ML venues}
\label{sec:lit-review}

Peer review at ML venues is noisy. When the same papers were reviewed by two separate committees, many papers accepted
by one were rejected by the other \citep{cortes2021inconsistency,beygelzimer2023has}. \citet{shah2022challenges}
surveys the problems of review at computer science venues, such as bias, miscalibration and strategic behaviour.

ICLR is a common object of study because its reviews and decisions are public on OpenReview. \citet{kang2018dataset}
released a dataset of drafts and reviews from ICLR and other venues and showed that acceptance is partly predictable
from the text. \citet{tran2020open} and \citet{zhang2022investigating} studied how acceptance at ICLR relates to authors
and institutions, and \citet{sun2022does} and \citet{manzoor2021uncovering} used ICLR's switch to double-blind review in
2018 to show that blinding reduced, but did not remove, a prestige effect. Where reviewers can see the authors, famous
authors fare better \citep{tomkins2017reviewer}, and preprints can reveal who wrote a paper even under double-blind
review \citep{bharadhwaj2020deanonymization}. \citet{ibrahim2026bias} tested a prestige gap at ICLR with a pre-registered
out-of-sample test. Other studies relate ICLR acceptance to the review process \citep{jung2025what} and to research
area \citep{wang2023what,xu2026reviewer}.

These studies ask mostly who gets accepted. They do not ask how acceptance relates to a paper's position in a changing
field: whether it is in a crowded topic, whether its authors are new to the topic, or how far it is from the team's
earlier work.

\section{Meta-research on machine learning and NLP}
\label{sec:lit-ml}

ML has a growing literature that examines its own research practice. \citet{lipton2019troubling} describe trends in
ML writing, such as claims that go beyond the evidence. Several authors argue that ML converges for reasons other than
merit: on ideas that suit the available hardware \citep{hooker2021hardware}, on a few benchmarks
\citep{dehghani2021benchmark,koch2021reduced,raji2021ai} and on foundation models that many applications share
\citep{bommasani2021opportunities}. Others note the growing role of industry and compute in deep learning
\citep{ahmed2020dedemocratization,abdalla2023elephant} and the values that highly cited ML papers express, among them
performance, novelty and building on past work \citep{birhane2022values}. On the other side, \citet{sutton2019bitter}
argues that convergence on general methods that scale with computation has been the efficient choice.

Studies of NLP describe paradigm shifts in the field. Interviews and citation analysis show cycles of exploration and
exploitation \citep{gururaja2023build,saphra2024first}, and an analysis of tasks, methods and datasets in NLP papers
found that methods and tasks lead and datasets follow \citep{pramanick2023diachronic}. \citet{klinger2020narrowing}
found that the thematic diversity of AI research has stagnated.

Most of this literature consists of essays, interviews or descriptive trends. Few studies test claims of herding or
homogenisation with hypotheses fixed in advance.

\section{Measuring research topics over time}
\label{sec:lit-topics}

Computational histories of a field usually fit a topic model to all its papers and then read how topics rise and fall.
\citet{griffiths2004finding} did this for PNAS abstracts, \citet{hall2008studying} and \citet{anderson2012towards} for
computational linguistics, and \citet{schafermeier2021topic} for ML. Dynamic topic models and BERTopic also use all
years when fitting \citep{blei2006dynamic,grootendorst2022bertopic}. \citet{prabhakaran2016predicting} showed that the
role a topic plays in abstracts, for example as a method or as a task, helps predict whether it will grow.

Judging which topics are emerging is hard, because emergence is partly defined by what happens later
\citep{rotolo2015what}, and a classification of topics fitted at one time may not suit later years
\citep{leydesdorff2002dynamic}. \citet{li2025evaluation} evaluated how well static topic models detect emerging topics,
and \citet{khanbayov2026topical} tested a signal for emerging AI topics that uses only earlier data. In prediction,
information that would not have been available at the time, known as leakage, makes results look better than they are
\citep{kaufman2012leakage,kapoor2023leakage}. Hindsight bias is the same problem in human judgement
\citep{fischhoff1975hindsight}.

This report also uses tools for confirmatory research on existing data. Pre-registration means fixing hypotheses and
analyses before seeing the outcomes, and it can be applied to data that already exist
\citep{nosek2018preregistration,vandenakker2021preregistration}. Multiverse and specification curve analyses run a
result under many reasonable choices to see how much it depends on them
\citep{steegen2016increasing,simonsohn2020specification}, and analyses chosen after looking at the data should be
treated as exploratory \citep{gelman2014statistical}. Labels produced by LLMs are increasingly used in such research,
but their agreement with human coders varies by task \citep{ziems2024can}.

\section{Summary and research gap}
\label{sec:lit-gap}

The literature offers two broad accounts of why a field converges: imitation of others, and continuity with one's own
record. It provides evidence on each, but mostly in separate settings and for published work. Studies of fashions
measure adoption rather than reward. Studies of novelty and newcomers measure reward through citations, grants or
journals and rarely see the decision on every submission. Studies of ML review focus on prestige and noise rather than
on herding, and histories of ML and NLP fit topics with hindsight.

This report brings these strands together in one population. It asks whether ICLR's waves survive measurement without
hindsight (RQ1), what kind of contribution newcomers bring to a topic (RQ2), and whether acceptance follows the crowd or
the team's own record (RQ3). It uses hypotheses registered in advance and tests the main ones again on held-out papers.
